% Transcription — fonds Alexandre Grothendieck, shelfmark 43, batch 5 (pages 81-100).
% Established from the facsimile archives/batches/43/batch-05.pdf (PDF page k =
% archivists' page 80+k, no cover sheet), cut from a copy of 43.pdf fetched
% through the project's own relay
% (https://grothendieck-relay.onrender.com/source/43.pdf); tiles in
% archives/tiles/43/p81-p100. Per the skill transcribe-grothendieck. The folder
% is 207 pages in eleven batches, transcribed in parallel; batches 1 and 9-11
% existed and were read for the folder's conventions (\underline for his
% underlined Hom/Ext/H, \mathbb{G}_m, \mathcal{C} for his curly C, apparatus
% style); batches 2-4 and 6-8 did not exist, so page 81 opens without the
% context of batch 4, and the notation here is fixed from these pages.
%
% Pass: Opus 5.5 (claude-opus-5-5), 2026-09-26 — first pass, unchecked against the pages by a human.
%
% Dating: the folder has its own inventory date, carried verbatim in
% \dating{} with no group sentence (the group « Jacobiennes » (43-44) is added
% by the skill only for undated folders). No page of this batch carries a date.
% The typed bibliography of pp. 84/86 cites Lang's Bourbaki talk of May 1964 and
% Néron's IHES paper of 1964 — a terminus for that typescript, not for his notes.
%
% Copies. No page of this batch can be identified with confidence as a
% photocopy. Pages 81, 83, 85 are ink (blue and black, with pencil frames on
% 83 and 85) on white paper — originals. Pages 88-89 and 91-97 are on grey
% paper with uniformly dark strokes and could be copies, but page 96 carries a
% blue ink trace over the grey, which points to an original; page 94 shows
% faint, right-reading traces of the text of page 93 (offset, not mirror
% show-through). Page 99 is pencil on yellowed paper, an original. All
% transcribed as originals; whoever checks against the leaves should settle
% 88-89 and 91-97 first.
%
% Pages skipped: 82, 87, 98 (blank cream sheets); 84 and 86 (two leaves of a
% typed bibliography, typed pages 131 and 132, [1]-[17] — Cartier, Lang,
% Lang-Néron, Néron, Northcott, Samuel, Serre, Shimura, Weil — someone
% else's thesis or memoir, with nothing in his hand); 100 (an upside-down
% typed page, « 22. », of an English paper on Thom complexes, Stiefel
% manifolds and James's « intrinsic join », Pembroke College, Cambridge —
% the back of p. 99, unrelated to the notes; it carries one short ink addition,
% « + up to homotopy type », in a hand not identified as his, keyed to the
% typed text). Page 90 is a cover (kraft folder) inscribed in his hand at the
% top right, « Corps de classe local géométrique et dualité »: it takes no
% \page marker and its words become the section heading, with a \note.
%
% Pages summarised: none. The large category diagrams of pp. 83 and 85 are
% redrawn as tikz-cd grids, their layout simplified and said in \note{}s; the
% oblique marginal notes of pp. 91, 93, 95 are given with a heavy apparatus.
%
% His pagination: none visible. Typed page numbers 131, 132 on the skipped
% bibliography; a pencilled « 99 » at the foot of p. 99 coincides with the
% archivists' number. Numbered runs of his: the four formulations (i)-(iv)
% of local duality on p. 88; the conditions (i)-(iii) on the prolonged
% Weil sheaf, p. 91; the a), b) of p. 93; and the work plan 1)-8) of p. 99
% (K-groupes algébriques; dualité de Cartier; cohomologies algébriques;
% dualité de Serre-Rosenlicht (?); le foncteur Gamma_{K/k}, « le Néron-... »;
% les foncteurs Ext*_{K/k}, H*_{K/k}; l'hom. fondamental; le th. de
% dualité), items 3), 5), 8) ticked. The run under the cover of p. 90
% (pp. 91-97) is continuous: p. 94 continues p. 93 (« Il faudrait aussi »),
% p. 96 opens with the cancelled tail of the N.B. at the foot of p. 95.
%
% What the batch is about. Duality for commutative algebraic groups over a
% perfect field and over a complete discrete valuation ring. P. 81: one line,
% Hom_{k-gr}(U_0, G) = Hom(J^0_{L/k} ⊗ E^1, G) (reading doubtful). P. 83:
% a diagram of the categories in play (k-group schemes, quasi-compact
% k-groups, locally algebraic preschemes, algebraic groups, pro-algebraic
% groups, fppf and fpqc sheaves), fppf sheaves as fpqc sheaves commuting
% with filtered limits of rings; F(S) = lim G_i(S) versus F'(S) = (lim G_i)(S),
% which coincide on the small fppf site, F' = u^*u_*(F) for u: fpqc -> fppf;
% the programme « compute the Ext^i of the elementary groups alpha_p, mu_p,
% Z/pZ, Z/lZ, G_a, G_m, Z, A over an algebraically closed k ». P. 85: the
% category (QSch) obtained by inverting integral radicial surjective
% morphisms, qfpqc and qfppf sheaves, and a lattice of the corresponding
% categories of groups « mod radiciels », with « k-gr. alg./radiciels ≃
% groupes quasi-algébriques ». P. 88: four formulations of local duality for a
% constructible complex F (bidual, H^i_a dual to H^{-i}, H^i(X - a) dual to
% H^{-i-1}, Ext with R i^!), each marked with the hypothesis it needs (« bon »,
% strictly local, henselian). P. 89: the same for U -> X and for a ∈ Z ⊂ Y ⊂ X,
% with two long exact sequences of local cohomology set face to face under
% the pairing. Pp. 91-97, under « Corps de classe local géométrique et
% dualité »: V a complete DVR with perfect residue field k, the ring scheme V
% and the ind-ring-scheme K = W(A)_pi; the « Weil topology » on S = Spec V
% (flat S' of finite type, fpqc descent), the category C(S) of abelian Weil
% sheaves and their canonical prolongation to flat preschemes; H^i, Ext^i,
% underlined Ext and the local-to-global spectral sequence; the generic
% open U, i_*: C(U) -> C(S) and « quels sont les foncteurs dérivés ?? »; the
% topology T'_k on k (descent by finite surjective morphisms), C'(k), H^i(k, F) = 0
% for k algebraically closed, a programme to compare with Galois cohomology;
% the « Néron » functor N: C(S) -> C'(k), N(F)(A) = lim_P prod F(V(A_i^P))
% over finite constructible partitions of A; N = H^0_{S/k}, its derived
% functors H^i_{S/k}, H^0_S = H^0_k ∘ H^0_{S/k}, the Leray spectral sequence
% H*(S, F) <= H^p(k, H^q_{S/k}(F)); relative Hom_{S/k}, Ext_{S/k} and the
% spectral sequences for Ext*(S; F, G), with « vérifier » in the margin. P. 99:
% Ext^i_{K-gr}(G, S gamma) ≅ Ext^{i+n}_{k-gr}(T^*G, gamma) for G affine
% without G_a component, and the eight-point plan above.
%
% Notation fixed once for the batch: \mathbb{V}, \mathbb{K}, \mathbb{W} for
% his double-struck V, K, W (p. 91); \mathcal{C} for his curly C (categories
% of sheaves); \mathcal{T} for his curly T (topologies); \mathcal{N} for the
% Néron functor; \mathcal{P} for his barred P (partitions, p. 95);
% \underline{\mathrm{Ext}}, \underline{\mathrm{Hom}}, \underline{H} for his
% underlined sheaf-valued functors (double underlining on p. 99 rendered
% single, said in a \note); \mathbb{E}\mathrm{xt} for his double-struck
% hyper-Ext (pp. 92-93); \hat{F} for the dual complex (pp. 88-89).
%
% Readings to check first: the only line of p. 81 (its subscript stack); the
% node labels of the diagrams of pp. 83 and 85; the oblique margins of
% pp. 91, 93, 95; the cancelled N.B. at the foot of p. 95 and head of p. 96;
% the name after « Serre- » and after « Néron- » on p. 99; the exponent 2i on
% p. 99.
% Apparatus: about 48 \ill{} and 190 \uncertain{} in the body.

\documentclass[11pt,a4paper]{article}
% Le préambule commun à toutes les transcriptions du fonds Grothendieck.
%
% Il tient en une idée : **l'incertitude se compose, elle ne se résout pas.**
% Une transcription qui lisse un mot douteux a détruit la seule chose pour
% laquelle on la faisait — un lecteur doit pouvoir savoir, à chaque endroit,
% ce qui était sur le feuillet et ce qui a été ajouté. Les six macros qui
% suivent sont donc l'appareil critique, et non de la décoration.
%
% Les mêmes commandes sont reconnues par `scripts/render.mjs`, qui produit la
% vue de lecture du site. Une macro ajoutée ici doit l'être aussi là-bas,
% faute de quoi le rendu échouera bruyamment — ce qui est voulu : un rendu
% silencieusement faux est pire qu'un rendu absent.

\ProvidesPackage{grothendieck}[2026/08/08 Transcriptions du fonds Grothendieck]

\usepackage{amsmath,amssymb,amsthm}
\usepackage{mathrsfs}
\usepackage{stmaryrd}
% Les diagrammes commutatifs sont partout dans ce fonds : c'est la langue
% même de la géométrie algébrique de Grothendieck, et une transcription qui
% les rendrait en prose ne transcrirait rien.
\usepackage{tikz-cd}
% Français par défaut — c'est la langue des feuillets — mais l'anglais est
% chargé aussi : la traduction et le résumé sont des documents anglais, et la
% typographie française (espaces avant les deux-points) y serait une faute.
% Ces documents basculent par \selectlanguage{english} en tête de corps.
\usepackage[english,french]{babel}
\usepackage{fontspec}
\usepackage{geometry}
\geometry{a4paper,margin=25mm,marginparwidth=28mm}
\usepackage{marginnote}
\usepackage{xcolor}
% ulem, not soul. `soul`'s \st and \ul are built on a character-by-character
% scan that XeLaTeX's font machinery breaks: the document fails with an
% undefined control sequence rather than degrading. ulem does the same two jobs
% and is Unicode-engine safe. `normalem` stops it hijacking \emph, which the
% transcriptions use for Grothendieck's own emphasis.
\usepackage[normalem]{ulem}
\usepackage{hyperref}
\hypersetup{colorlinks=true,linkcolor=black,urlcolor=black,pdfborder={0 0 0}}

% --- Métadonnées du lot -----------------------------------------------------
% Reprises telles quelles par le script de rendu, qui les affiche en tête de
% la vue de lecture. Elles fixent ce que le fichier prétend couvrir : sans
% elles, une transcription flotte sans point d'attache dans un fonds de seize
% mille feuillets.

\newcommand{\folder}[1]{\gdef\grothFolder{#1}}
\newcommand{\batch}[1]{\gdef\grothBatch{#1}}
\newcommand{\leaves}[2]{\gdef\grothFirst{#1}\gdef\grothLast{#2}}
\newcommand{\foldertitle}[1]{\gdef\grothFoldertitle{#1}}
\gdef\grothFolder{?}\gdef\grothBatch{?}\gdef\grothFirst{?}\gdef\grothLast{?}\gdef\grothFoldertitle{}

% --- Le résumé --------------------------------------------------------------
%
% Il ouvre la lecture modernisée, sous le titre « Résumé », et il est composé
% en petit. Ce n'est pas un ornement : c'est le seul passage du document qui
% ne soit pas des mathématiques, et un lecteur doit pouvoir le reconnaître
% avant de l'avoir lu — puis le sauter s'il connaît déjà le sujet, ou s'y
% arrêter s'il ne le connaît pas. Le filet qui le suit marque la reprise du
% corps.

\newenvironment{resume}
  {\small\setlength{\parskip}{0.45em}}
  {\par\medskip\hrule\medskip}

% --- Mots-clés ---------------------------------------------------------------
%
% En anglais, en clôture du résumé de la lecture modernisée : le vocabulaire
% moderne sous lequel chercher ce que les feuillets construisent. Le manifeste
% du site (`npm run manifest`) les extrait pour étiqueter la cote entière —
% cette ligne est la source unique des tags, il n'y a pas de fichier à côté.
\newcommand{\keywords}[1]{\par\smallskip\noindent
  {\footnotesize\textsc{Keywords} --- \textit{#1}}\par}

% --- L'appareil critique ----------------------------------------------------

% Le numéro de feuillet, en marge. C'est l'ancre qui permet de régler un
% désaccord sur une formule en regardant *un* feuillet plutôt qu'une chemise
% de six cents. Le site s'en sert aussi pour tourner les pages du fac-similé
% quand on fait défiler la transcription.
\newcommand{\leaf}[1]{%
  \marginnote{\footnotesize\color{gray!70}\textsf{#1}}%
  \label{leaf:#1}\ignorespaces}

% Illisible : on ne devine pas. Un mot inventé qui se lit comme les autres est
% le pire résultat possible.
\newcommand{\ill}{\textcolor{red!60!black}{[\,\ldots\,]}}

% Lecture douteuse : le mot est donné, et signalé comme incertain.
\newcommand{\uncertain}[1]{\uline{#1}}

% Ajout de l'éditeur — une lettre manquante, un mot sous-entendu.
\newcommand{\add}[1]{\textcolor{blue!50!black}{[#1]}}

% Biffé par Grothendieck lui-même. Ce qu'il a rayé fait partie du document :
% une rature dit souvent où la pensée a changé de direction.
\newcommand{\struck}[1]{\sout{#1}}

% Note du transcripteur, sur l'état matériel du feuillet.
\newcommand{\note}[1]{\par\smallskip{\small\color{gray!60!black}\textit{[#1]}}\par\smallskip}

% Note marginale de Grothendieck, distincte du corps du texte.
\newcommand{\marginal}[1]{\par\smallskip\noindent\hspace{1em}%
  {\small\color{gray!60!black}#1}\par\smallskip}

% --- Titre ------------------------------------------------------------------

\AtBeginDocument{%
  \begin{center}
    {\large\bfseries Fonds Alexandre Grothendieck --- cote n\textsuperscript{o}~\grothFolder}\\[2pt]
    {\itshape\grothFoldertitle}\\[4pt]
    {\small feuillets \grothFirst--\grothLast\ (lot \grothBatch)}\\[2pt]
    {\footnotesize\color{gray!60!black}Transcription établie d'après le fac-similé numérisé
     par l'Université de Montpellier.\\
     Les lectures incertaines sont soulignées, les illisibles notées [\,\ldots\,],
     les ajouts entre crochets.}
  \end{center}
  \vspace{1em}}

\endinput


\folder{43}
\batch{5}
\pages{81}{100}
\dating{[à partir de 1961-vers 1968]}
\watermark{Édition de démonstration}
\foldertitle{Dualité des groupes. Jacobiennes généralisées etc : notes et copies de notes manuscrites (s.d.)}

\begin{document}

\page{81}
\[
\mathrm{Hom}_{k\text{-}\mathrm{gr}}(U_0, G) = \mathrm{Hom}_{\ill{}}\bigl(J^{0}_{L/k} \otimes E^{1}, G\bigr)
\]
\note{la seule ligne du feuillet, en haut. L'indice du second $\mathrm{Hom}$
est un empilement de petites lignes serrées : au-dessus d'un trait,
\uncertain{« vu limite »} ; au-dessous, \uncertain{« multiplicatif »}, puis
\uncertain{« pour un pro- »} \ill{} \uncertain{« de groupes »}. Le signe entre
$J^{0}_{L/k}$ et $E^{1}$ est lu $\otimes$ sans certitude ; le $U_0$ est une
lecture de son $\mathcal{U}$ bouclé.}

\page{83}
\note{la moitié supérieure du feuillet est un schéma d'inclusions entre
catégories, tracé à l'encre avec, au crayon, de grands traits obliques
qui délimitent des régions ; il est redessiné ici en grille, la disposition
simplifiée. Les flèches sont celles de la page ; les libellés des nœuds sont
abrégés comme il les abrège.}

\begin{tikzcd}[column sep=small, row sep=small, nodes={font=\scriptsize}]
 & \text{faisceaux fpqc} & \\
\text{$k$-schémas en groupes} \arrow[ur] & & \text{faisceaux fppf} \arrow[ul] \\
\text{$k$-gr.\ qu.-cpts} \arrow[u] \arrow[drr, dashed] & \text{$k$-préschémas loc.\ alg.} \arrow[ul] \arrow[ur] & \text{gr.\ pro-alg.} \arrow[u] \\
 & \text{$k$-gr.\ alg.} \arrow[ul, "F"'] \arrow[u] \arrow[r] & \text{$k$-gr.\ qu.-cpts $F_s$} \arrow[u]
\end{tikzcd}

\note{sur la flèche de « faisceaux fppf » vers « faisceaux fpqc », une petite
croix. Dans les nœuds, « qu.-cpts » (deux fois) et « préschémas » sont des lectures
douteuses ; le nœud « $k$-préschémas loc.\ alg. » a un premier mot surchargé. Sur la flèche de « $k$-gr.\ \uncertain{qu.-cpts} $F_s$ » vers
« gr.\ pro-alg. », un petit encadré : « \uncertain{définissable} \uncertain{fl.\
fid.}\ ? ». À gauche, entre parenthèses, « (ne sont pas abéliennes) », relié
par des flèches pointillées à « $k$-schémas en groupes » et au nœud central.}

À droite du schéma, encadré :
\begin{quote}
[faisceaux fppf] $=$ faisceaux fpqc qui commutent aux $\varinjlim$ d'anneaux
(filtrantes) \uncertain{?}
\end{quote}
et, en face de « gr.\ pro-alg. » : « \ill{} \ill{} \emph{surjectives} (pas
injectives !) » ; en face de « $k$-gr.\ \uncertain{qu.-cpts} $F_s$ »,
encadré : « $\simeq$ gr.\ proalg.\ \uncertain{con.}\ affines ? ».

$(G_i)_{i \in I}$ \uncertain{système} \uncertain{direct} \uncertain{de}
faisceaux fpqc.
\begin{align*}
&(\mathrm{i}) \quad F(S) = \varinjlim_i G_i(S) \\
&(\mathrm{ii}) \quad F'(S) = \Bigl(\varprojlim G_i\Bigr)(S)
\end{align*}
\note{sous la limite projective de (ii) : « (\uncertain{dans} \uncertain{cat.}\
des préschémas) ».}
\marginal{Coïncident sur le \emph{petit} site fppf}
\[
\boxed{F' = u^{*} u_{*}(F)}
\]
\begin{tikzcd}
\text{site fpqc} \arrow[d, "u"] \\
\text{site fppf}
\end{tikzcd}

$u^{*} = $ \uncertain{restriction} \uncertain{évidente} des sites / $u^{*}$
\uncertain{défini} par image inverse \uncertain{usuelle} \uncertain{pr.}\
\uncertain{usuel}

$u_{*} = $ restriction des foncteurs
\note{un « $u$ » isolé suit, sans suite. La fin de la ligne de $u^{*}$,
serrée en bout de ligne sur deux étages, est très douteuse.}

\begin{quote}
Calculer les $\mathrm{Ext}^{i}$ des groupes « élémentaires » sur \uncertain{corps}
alg.\ clos $k$, en vérifiant que c'est \uncertain{pareil} pour les catégories
envisagées (pour autant que ça ait un sens\ldots)
\[
\alpha_p \quad \mu_p \quad \mathbb{Z}/p\mathbb{Z} \quad \mathbb{Z}/\ell\mathbb{Z} \quad \mathbb{G}_a,\ \mathbb{G}_m,\ \mathbb{Z},\ \mathbf{A}
\]
\end{quote}
\note{une accolade sous $\mu_p$, $\mathbb{Z}/p\mathbb{Z}$ et une autre sous
$\mathbb{G}_m$, $\mathbb{Z}$ ; le $\mathbb{Z}$ est surchargé ; $\mathbf{A}$ en
très grande capitale, sans doute la variété abélienne.}
\begin{quote}
Faire le calcul dans les catégories \uncertain{correspondantes} \uncertain{en se}
limitant à ce qui est \uncertain{annulé} \uncertain{par} \uncertain{un}
\uncertain{$n$} fixé.
\end{quote}

\page{85}
(\uncertain{QSch}) \uncertain{Catégorie} \uncertain{obtenue} à partir de
\uncertain{(Sch)} \uncertain{en} \uncertain{rendant} \uncertain{inversibles}
les morphismes entiers radiciels surjectifs.

\uncertain{Parallèlement}, fppf et fpqc définis ; on parle de faisceaux qfpqc
et qfppf, \uncertain{ce} \uncertain{sont} les faisceaux fpqc et fppf qui
transforment morphismes entiers-surj.\ rad.\ en isom.

\note{au-dessous, un treillis de catégories aux traits sans pointe (sauf les
deux flèches vers le nœud du haut), dans trois cadres au crayon ; il est
redessiné en grille, la disposition simplifiée. Les libellés sont abrégés
comme il les abrège.}

\begin{tikzcd}[column sep=small, row sep=small]
 & & (1) & \\
(2) \arrow[urr] \arrow[d, no head] & & & \\
(3) \arrow[dr, no head] & (4) \arrow[uur] \arrow[ul, no head] \arrow[d, no head] \arrow[dr, no head] & & (5) \arrow[d, no head] \\
 & (6) \arrow[dr, no head] & (7) \arrow[ur, no head] \arrow[d, no head] & (8) \\
 & & (9) \arrow[r, no head] & (10) \arrow[u, no head]
\end{tikzcd}

\begin{itemize}
\item[(1)] faisceaux sur le site des $k$-schémas qu.-cpts parfaits
\item[(2)] groupes de la cat.\ des $k$-schémas loc.\ qu.-cpts parfaits
\item[(3)] groupes de la catégorie des $k$-schémas qu.-cpts parfaits
\item[(4)] $k$-sch.\ en gr.\,/\,rad.
\item[(5)] faisceaux sur le site des $k$-schémas loc.\ qu.-cpts parfaits essentiellement de t.f.\ sur $k$
\item[(6)] $k$-gr.\ qu.-cpts, mod rad
\item[(7)] $k$-gr.\ loc.\ alg.\,/\,radiciels
\item[(8)] gr.\ pro-quasi-alg.
\item[(9)] $k$-gr.\ alg.\,/\,radiciels $\simeq$ groupes quasi-alg.\ \uncertain{pas} de t.f.
\item[(10)] $k$-gr.\ qu.-cpts\,/\,radiciels
\end{itemize}
\note{les numéros (1)--(10) sont de l'éditeur : les libellés, trop longs pour
tenir dans les nœuds, sont reportés dans la liste.}

\note{dans (2), « qu.-cpts » est surchargé en « loc.\ qu.-cpts » ; au-dessus de (4), un « groupes de la cat.\ » est biffé et
récrit. Sous (7), une ligne biffée
(« \struck{\ill{} préschémas loc.\ \ill{}} ») ; dans (5), un
premier « qfppf » est biffé. Sous (10), un « mod rad » biffé. Les deux traits qui aboutissent à (1) portent
une pointe ; tous les autres n'en ont pas. Dans (9), le mot entre « quasi-alg. »
et « de t.f. » est très incertain.}

\page{88}
Formulations de la \emph{dualité locale} ($F^{\bullet}$ constr.)
\note{quatre énoncés numérotés (i)--(iv), reliés dans la marge gauche par des
doubles flèches verticales : (i) $\Downarrow$ (ii), (ii) $\Updownarrow$ (iii),
(iii) $\Updownarrow$ (iv), et une longue flèche hachurée de (i) à (iii). En
face, à droite, une colonne de croix et de conditions.}

(i) $F^{\bullet} \to D D F^{\bullet}$ \uncertain{est} \uncertain{un}
\uncertain{isom.}\ \struck{pour $F^{\bullet}$ constructible \ill{}}
— $\times$ « bon »

(ii) a) $H^{i}_{a}(X, F^{\bullet})$ dual de $H^{-i}(X, \hat{F}^{\bullet})$
— $\times$ \emph{strictement local} et « bon »

\struck{b) $H^{i}_{a}(X, F^{\bullet}) =$ \ill{} \ill{} \ill{}} $\Rightarrow$
\struck{$H^{i}(X, F^{\bullet}) = \mathcal{H}^{i}(F^{\bullet})_{x}$ \ill{}}
\note{la ligne b) est barrée d'un grand trait ondulé ; ce qui suit
« $\Rightarrow$ » est surchargé et biffé à son tour.}

(iii) $H^{i}(X - a, F^{\bullet})$ dual de $H^{-i-1}(X - a, \hat{F}^{\bullet})$
— \uncertain{id.}

(iv) $\mathrm{Ext}^{i}(X; F, \lambda(G)) \simeq \mathrm{Ext}^{i}(k; \mathbb{R}i^{!}(F), G)$
— $\times$ \uncertain{hensélien} « bon », corps \uncertain{résiduel}
\uncertain{$k$}
\note{au-dessus de « hensélien », « local » entouré. Le $a$ de (ii) et (iii)
est le point fermé (« $X - a$ ») ; le chapeau $\hat{F}^{\bullet}$ désigne le
dual. Le « $\mathrm{Ext}^{i}$ » du second membre de (iv) est très cursif.}
\note{sous un trait horizontal, le reste du feuillet est vide, à l'exception
d'une paire de parenthèses vides.}

\page{89}
\[
H^{i}(U, F) = H^{i}\bigl(X, \mathbb{R}i_{*}(F|U)\bigr) \quad \text{dual de}
\]
\[
U \xrightarrow{\ i\ } X, \qquad U \xrightarrow{\ i'\ } (X - a) \qquad
H^{-i}_{a}\bigl(X, \mathbb{R}i_{!}(\hat{F}|U)\bigr) \simeq H^{-i-1}\bigl(X - a, \mathbb{R}i'_{!}(\hat{F}|U)\bigr)
\]
\note{puis un « $H^{-}$ » sans suite, et plus bas un « $H^{i}$ » biffé.}

$a \in Z \subset Y \subset X$
\[
Y - Z \xrightarrow{\ i'\ } X - Z, \qquad Y - Z \xrightarrow{\ k\ } Y - a, \qquad
Y \xrightarrow{\ i\ } X, \qquad Z \xrightarrow{\ j\ } X
\]
\[
\begin{cases}
H^{i}_{Y-Z}(F) = H^{i}\bigl(Y - Z, \mathbb{R}i'^{!}(F|X - Z)\bigr) \\
\quad \text{dual de } H^{-i-1}\bigl(Y - a, \mathbb{R}k_{!}(\hat{F}|Y - Z)\bigr)
\end{cases}
\]
\note{dans le premier membre, un $H^{i}$ est surchargé en $H^{i}(Y-Z, \ldots)$ ;
dans le dernier, un « $\mathbb{R}k_{!}$ » est très surchargé, la lecture en est
douteuse.}
\[
\begin{cases}
H^{i}_{Y}(F) = H^{i}\bigl(Y, \mathbb{R}i^{!}(F)\bigr) \quad \text{dual de} \\
H^{-i}_{a}\bigl(Y, \mathbb{L}i^{*}(\hat{F})\bigr) = H^{-i}_{a}(Y, \hat{F}|Y)
\end{cases}
\]
\note{dans les deux lignes, un $X$ est biffé et remplacé par $Y$ ; le second membre de la dernière ligne était
d'abord précédé d'un signe biffé.}
\[
H^{i}_{Z}(F) \quad \text{dual de} \quad H^{-i}_{a}(Z, \hat{F}|Z)
\]
\[
\begin{array}{ccc}
H^{i}_{Z}(F) & \times & H^{-i}_{a}(\hat{F}|Z) \\
\downarrow & & \uparrow \\
H^{i}_{Y}(F) & \times & H^{-i}_{a}(\hat{F}|Y) \\
\downarrow & & \uparrow \\
H^{i}_{Y-Z}(F) & \times & H^{-i-1}\bigl(Y - a, \mathbb{R}k_{!}(\hat{F}|Y - Z)\bigr) \\
\downarrow & & \uparrow \\
H^{i+1}_{Z}(F) & \times & H^{-i-1}_{a}(\hat{F}|Z) \\
\downarrow & & \uparrow \\
H^{i+1}_{Y}(F) & \times & H^{-i-1}_{a}(\hat{F}|Y)
\end{array}
\]
\note{sur la page, les deux suites exactes longues sont écrites en deux lignes
horizontales mises en face, la seconde de droite à gauche, des croix marquant
les accouplements terme à terme ; on les dispose ici en colonnes. L'indice $a$ du premier terme de la seconde
ligne est surchargé. Le dernier terme, serré au bord, est lu
$H^{-i-1}_{a}(\hat{F}|Y)$ sans certitude.}

\section*{Corps de classe local géométrique et dualité}
\note{titre inscrit de sa main en haut à droite d'une chemise de papier
kraft (page 90), « Corps de classe local » souligné, sans autre contenu ;
il ouvre la suite, pages 91 et suivantes.}

\page{91}
$\mathbb{V}$ anneau de valuation discrète, complète, à corps résiduel $k$
parfait, corps des fractions $K$, uniformisante $\pi$.

$\mathbb{V}$ le schéma en anneaux sur $k$ défini par $\mathbb{V}$
\note{les deux $\mathbb{V}$ sont écrits à double trait, l'un gras, l'autre
maigre ; de même $\mathbb{K}$ et $\mathbb{W}$ ci-dessous.}

$\mathbb{K}$ le ind-schéma en anneaux défini par
\[
A \rightsquigarrow \mathbb{W}(A)_{\pi} = \mathbb{K}(A)
\]
\note{au-dessus de la ligne suivante, un $\mathcal{C}_S$ isolé.}

$S = \operatorname{Spec}(\mathbb{V})$, on munit $S$ de la « topologie de Weil »
définie par
\[
\left\lbrace
\begin{array}{l}
\text{catégorie des $S'$ plats de type fini sur $S$.} \\
\text{descente fid.\ plate [\uncertain{ou} fid.\ \uncertain{plate} \ill{} ???]}
\end{array}\right.
\]

\marginal{N.B. \uncertain{y a-t-il} \uncertain{lieu} \uncertain{de}
\uncertain{prendre} \uncertain{$S'$} \uncertain{simple} sur $S$ ? \uncertain{Non},
\uncertain{car} on \uncertain{ne} \uncertain{peut} \uncertain{pas}
\uncertain{prendre} \uncertain{un} \uncertain{gr.} \ill{} \uncertain{si}
$S_0'$ \uncertain{n'est} \uncertain{pas} \struck{\ill{}} \uncertain{réduit} ??
\uncertain{Peut-être}\ldots}
\note{note oblique dans la marge gauche, à l'encre plus pâle ; la lecture en
est presque entièrement douteuse. Dans l'accolade, la fin de la seconde ligne
est soulignée d'un trait ondulé.}

$\mathcal{C}(S) = $ catégorie des faisceaux de Weil abéliens sur $S$

Donc un $G \in \mathcal{C}(S)$ est un foncteur $S' \mapsto G(S')$, soumis
à des conditions d'exactitude connues.

On \uncertain{convient} de \emph{prolonger} $G$ en $\overline{G}$ aux
préschémas plats $S'$-\uncertain{préschémas} sur $S$, par les conditions
suivantes
\begin{itemize}
\item[(i)] $\overline{G}$ de nature locale
\item[(ii)] $\overline{G}$ commute aux limites inductives d'anneaux
\end{itemize}
Alors \ill{} \uncertain{aussi} \uncertain{(iii)}
\begin{itemize}
\item[(iii)] $\overline{G}$ compatible avec descente par morphismes fid.\
plats de présentation finie.
\end{itemize}

Exemple : tout préschéma en gr.\ \struck{sur} \uncertain{commutatif} $G$ sur
$S$ définit un faisceau de Weil, désigné par la même lettre. Si $G$ est loc.\
de type fini sur $S$, alors le prolongement \uncertain{canonique}
$\overline{G}$ est encore « représenté » par le \uncertain{même} $G$,
\uncertain{i.e.}
\[
\overline{G}(S') = \mathrm{Hom}_{S}(S', G).
\]
\note{« commutatif $G$ » est ajouté en interligne au-dessus d'un mot biffé.}

\page{92}
Dans la catégorie $\mathcal{C}(S)$, \uncertain{on} \uncertain{a} les foncteurs
$H^{i}(S, F)$ et $\mathrm{Ext}^{i}(S; F, G)$, \struck{\ill{}}
$\underline{\mathrm{Ext}}^{i}_{S}(F, G)$. \uncertain{Notations}, \uncertain{à}
\uncertain{ce} propos
\[
\begin{cases}
H^{i}(S, F) = \mathrm{Ext}^{i}(S; \mathbb{Z}_{S}, F) \\[4pt]
\mathbb{E}\mathrm{xt}^{*}(S; F, G) \Leftarrow H^{p}\bigl(S, \underline{\mathrm{Ext}}^{q}(F, G)\bigr)
\end{cases}
\]
\note{le second $\mathrm{Ext}$ de l'accolade est écrit avec un E à double
trait ; la flèche $\Leftarrow$ est celle d'une suite spectrale. L'indice de
$\mathrm{Ext}^{i}(S; F, G)$ à la première ligne est surchargé.}

Soit $U$ l'ouvert générique de $S$, d'où une catégorie de faisceaux
$\mathcal{C}(U)$. On \uncertain{a} un foncteur image directe par
$i \colon U \to S$
\[
i_{*} \colon \mathcal{C}(U) \to \mathcal{C}(S)
\]
Quels sont les foncteurs dérivés ??

\page{93}
Soit $\mathcal{T}'_{k}$ la \uncertain{topologie} de Weil \uncertain{suivante} :
\uncertain{on} \uncertain{adopte} \uncertain{sur} les \struck{\uncertain{quasi}
\uncertain{quasi-algébriques}} \struck{\ill{}} préschémas \uncertain{fl.}\
\uncertain{qu.-cpts}, la descente \uncertain{est} donnée par les morphismes
\struck{\ill{}} \emph{finis} surjectifs. Soit $\mathcal{C}'(k)$ la catégorie
des faisceaux de Weil sur $\mathcal{T}'_{k}$.
\marginal{on \uncertain{devrait} \uncertain{de} « finis » \uncertain{sans}
\uncertain{doute} \uncertain{pour} \uncertain{ce} \uncertain{dessin}
\uncertain{d'avoir} « application » \uncertain{entre} $\mathcal{T}_{S}$
\uncertain{et} $\mathcal{T}'_{k}$ !}
\note{note oblique dans la marge gauche, cernée d'un trait courbe ; lecture
presque entièrement douteuse. Le mot biffé devant « finis » est long et
surchargé ; « finis » est souligné deux fois.}

On \uncertain{a} dans $\mathcal{C}'(k)$ les foncteurs $H^{i}(k, G)$,
$\mathrm{Ext}^{i}(k; F, G)$, \uncertain{et} $\underline{\mathrm{Ext}}^{i}_{k}(F, G)$
\uncertain{et} les relations habituelles
\[
\begin{cases}
H^{i}(k, F) \simeq \mathrm{Ext}^{i}(k; \mathbb{Z}_{k}, F) \\[4pt]
\mathbb{E}\mathrm{xt}^{*}(k; F, G) \Leftarrow H^{p}\bigl(k, \underline{\mathrm{Ext}}^{q}_{k}(F, G)\bigr)
\end{cases}
\]
\note{« et $\underline{\mathrm{Ext}}^{i}_{k}(F, G)$ » est ajouté en interligne ;
dans la seconde ligne, l'argument $k$ est surchargé et les indices du dernier
$\underline{\mathrm{Ext}}$ sont serrés.}

On \uncertain{montre} que si $k$ est alg.\ clos, alors $H^{i}(k, F) = 0$ pour
$i > 0$. La formation des $\underline{\mathrm{Ext}}^{i}_{k}(F, G)$ commute :
l'extension du corps de base\ldots.

Il y a lieu de \struck{\uncertain{calculer}} \struck{\uncertain{donner}} des
règles de calcul pour des $H^{i}(k, G)$, $\mathrm{Ext}^{i}(k; F, G)$ si $G$
provient d'un \uncertain{quasi-}préschéma en groupes habituel sur $k$,
[p.\ ex.\ d'un pro-groupe de Serre de type « spécial », i.e.\ à morphismes
de transition \uncertain{affines} pour l'indice grand\ldots].
\note{« quasi » est ajouté en interligne au-dessus de « préschéma ».}
Il faudrait prouver alors que (\uncertain{k.l.} \ill{} \uncertain{p.\ ex.})
\begin{itemize}
\item[a)] si $G$ est simple sur $k$, ses $H^{i}(k, G)$ sont les
$H^{i}(k, \tilde{G})$, où $\tilde{G}$ est le faisceau de Weil associé\ldots
\item[b)] \emph{Kif-kif} pour les $\mathrm{Ext}^{i}(k; F, G)$\ldots.
\end{itemize}

\page{94}
Il faudrait aussi, de façon générale [\uncertain{dans} le \uncertain{contexte}
« quasi »] étudier \struck{\ill{}} si $H^{i}(k, G)$ peut se calculer comme une
cohomologie galoisienne\ldots.
\note{trois traits verticaux dans la marge gauche, en face de ce paragraphe.}

\struck{N.B. Il \uncertain{suffit} \uncertain{par} \uncertain{cohérence}
\uncertain{fid.}\ \uncertain{plates} \ill{} \ill{} Weil \uncertain{permutés}}
\note{ligne biffée d'un double trait et encadrée ; lecture très douteuse.
Le reste du feuillet est blanc.}

\page{95}
\marginal{(dit de Néron)}
On va définir un foncteur fondamental
\[
\mathcal{N} \colon \mathcal{C}(S) \to \mathcal{C}'(k)
\]
\note{« (dit de Néron) », « Néron » souligné, est écrit en haut à droite dans
une accolade qui pointe vers « fondamental ».}
Pour ceci, si $F$ est un \struck{faisceau} faisceau de Weil sur $S$, on définit
un faisceau de Weil sur $\mathcal{T}_{k}$ en prenant, pour $A$
\uncertain{une} \uncertain{algèbre} \uncertain{réduite} \uncertain{sur}
\uncertain{$k$}\ldots\ \ill{}.
\[
\mathcal{N}(F)(A) = \varinjlim_{\mathcal{P}} F\Bigl(\operatorname{Spec}\bigl(\mathbb{V}\bigl(\textstyle\prod_{i \in I_{\mathcal{P}}} A^{\mathcal{P}}_{i}\bigr)\bigr)\Bigr)
= \varinjlim_{\mathcal{P}} \prod_{i \in I_{\mathcal{P}}} F\bigl(\mathbb{V}(A^{\mathcal{P}}_{i})\bigr)
\]
\note{au-dessus de la première égalité, une variante surchargée,
qui s'écrit
« $= F\bigl(\varinjlim_{\mathcal{P}} \mathbb{V}(A^{\mathcal{P}})\bigr)$
$= F\bigl(\varprojlim_{\mathcal{P}} \operatorname{Spec}(\mathbb{V}(A^{\mathcal{P}}))\bigr)$ »,
avec, en exposant, « $\operatorname{Spec}(\varinjlim_{\mathcal{P}} \mathbb{V}(A^{\mathcal{P}}))$ »
relié par un « $=$ ». Sous le produit, une accolade donne
« $\prod_{i \in I_{\mathcal{P}}} A^{\mathcal{P}}_{i} = A^{\mathcal{P}}$ ». Le
$\mathcal{P}$ transcrit ici est chez lui un P barré (une partition) ; le « Spec »
est surchargé sur un S.}

$\mathcal{P}$ parcourant les « partitions constructibles \uncertain{finies} »
de $A$, et où pour une telle $\mathcal{P}$, $I_{\mathcal{P}}$ est son ensemble
d'indices, et $A^{\mathcal{P}}_{i}$ les composantes correspondantes de $A$.
\uncertain{En} d'autres termes, on \uncertain{prend} \uncertain{un}
\uncertain{préschéma} fl.\ \uncertain{qu.-cpt} \uncertain{réduit}
(\uncertain{cqch}) \ill{}, considéré comme un pro-(schéma affine) sur $k$,
\uncertain{on} le transforme par $\mathbb{V}$ en un pro-(schéma affine) sur $S$,
et on lui applique le foncteur $F$.
\struck{[N.B. $\mathcal{N}(F)$ \uncertain{ne} \uncertain{commute}
\uncertain{pas} \ill{} \uncertain{les} \uncertain{limites} \ill{} \ill{}
\ill{}\ldots\ \uncertain{pratiquement}, \uncertain{on} \uncertain{pourrait}
\uncertain{se} \uncertain{limiter}}
\note{le « [N.B.\ldots » du bas est biffé de grandes croix ; lecture très
douteuse. « cqch » (quelque chose) est ajouté en interligne.}

\marginal{N.B. On a une « application » \ill{} \ill{} \ill{}
$\mathcal{C}'_{k} \to \mathcal{C}_{S}$, \uncertain{produits}, \uncertain{dit}
\uncertain{le} foncteur \uncertain{\ill{}} \uncertain{faisceaux}
\uncertain{formels} \ill{}\ldots\ \ill{} \ill{} \uncertain{produits}
\uncertain{directs} ! ?}
\marginal{Non, car le foncteur $\mathbb{V} \colon k\text{-alg.} \to
\mathbb{V}\text{-alg.}$ \uncertain{ne} \uncertain{commute} \uncertain{pas}
\uncertain{aux} limites inductives\ldots}
\note{la première note marginale est oblique, cernée d'un trait courbe dans
la marge gauche ; la seconde, en bas à gauche, également oblique ; toutes
deux sont d'une lecture presque entièrement douteuse. Le $\mathcal{C}'_{k}$
de la première est surchargé.}

\page{96}
\struck{dans la catégorie $\mathcal{C}'(k)$ \uncertain{un} foncteur
\uncertain{commutant} \uncertain{aux} \ill{} \uncertain{avec} les limites
inductives}
\note{ces deux lignes, suite de la note biffée du bas de la page 95, sont
barrées de grandes croix.}

On constate que le foncteur $\mathcal{N}$ \uncertain{qu'on} vient de définir
est exact à gauche, \uncertain{donc} transforme groupe injectif en groupe
injectif\ldots, d'où
$\mathcal{C}(S) \to \mathcal{C}'(k)$. On \uncertain{pose} \uncertain{aussi}
\[
\mathcal{N}(F) = \underline{H}^{0}_{S/k}(F)
\]
et \struck{\ill{}} par suite
\[
\mathrm{R}^{i}\mathcal{N}(F) = \underline{H}^{i}_{S/k}(F)
\]
Notons l'isomorphisme
\[
H^{0}_{S} \simeq H^{0}_{k} \circ \underline{H}^{0}_{S/k}
\]
et d'autre part que $\underline{H}^{0}_{S/k}$ \emph{transforme injectifs en
injectifs}
\struck{\uncertain{car} \uncertain{construit} \ill{} \uncertain{adjoint} \ill{}
$\mathcal{N} = \underline{H}^{0}_{S/k}$ \ill{} \uncertain{par} \ill{}
\uncertain{dernier} \uncertain{est} \uncertain{exact} \ill{} \uncertain{généralisé} :
$\mathrm{Hom}(\Phi, \mathcal{N}(F)) \simeq \mathrm{Hom}(\mathcal{N}^{*}(\Phi), F)$}
\marginal{à vérifier\ldots}
\note{« transforme injectifs en injectifs » est souligné ; « à vérifier\ldots »
est écrit dans la marge gauche en face. Les trois lignes qui suivent, avec la
formule d'adjonction, sont encadrées et biffées de grands traits ; le
$\mathcal{N}^{*}$ de la formule est une lecture douteuse.}
D'où une suite spectrale de Leray
\[
H^{*}(S, F) \Leftarrow H^{p}\bigl(k, \underline{H}^{q}_{S/k}(F)\bigr)
\]

\page{97}
Plus généralement, définissons
\[
\underline{\mathrm{Hom}}_{S/k}(F, G) = \underline{H}^{0}_{S/k}\bigl(\underline{\mathrm{Hom}}_{S}(F, G)\bigr)
\]
\struck{Alors} et \uncertain{notons} les foncteurs dérivés de
$G \mapsto \underline{\mathrm{Hom}}_{S/k}(F, G)$ par
\[
\underline{\mathrm{Ext}}^{i}_{S/k}(F, G)
\]
On a alors, si on admet que pour $G$ injectif,
$\underline{\mathrm{Hom}}_{S}(F, G)$ est acyclique pour
$\underline{H}^{0}_{S/k}$ \struck{\uncertain{et} \uncertain{pour} $H^{0}_{S}$},
la suite spectrale
\[
\underline{\mathrm{Ext}}^{*}_{S/k}(F, G) \Leftarrow \underline{H}^{p}_{S/k}\bigl(\underline{\mathrm{Ext}}^{q}_{S}(F, G)\bigr)
\]
\marginal{vérifier}
\note{un trait vertical dans la marge gauche longe l'hypothèse, avec
« vérifier » en face ; l'indice de $\underline{\mathrm{Ext}}^{*}$ est
surchargé d'une tache d'encre.}
et compte tenu de
\[
\mathrm{Hom}_{S}(F, G) \simeq H^{0}\bigl(k, \underline{\mathrm{Hom}}_{S/k}(F, G)\bigr)
\]
et admettant que pour $G$ injectif, $\underline{\mathrm{Hom}}_{S/k}(F, G)$ est
acyclique pour $H^{0}(k, -)$, on trouve
\[
\mathrm{Ext}^{*}(S; F, G) \Leftarrow H^{p}\bigl(k, \underline{\mathrm{Ext}}^{q}_{S/k}(F, G)\bigr)
\]
\marginal{vérifier}
\note{même trait vertical et même « vérifier » en marge.}
\uncertain{Signalons} \uncertain{aussi}
\[
\underline{H}^{i}_{S/k}(F) \simeq \underline{\mathrm{Ext}}^{i}_{S/k}(\mathbb{Z}_{S}, F)
\]

Notations et remarques analogues \uncertain{en} \uncertain{partant} de
faisceaux sur $U \subset S$, au lieu de $S$.

\page{99}
\[
\underline{\mathrm{Ext}}^{i}_{K\text{-}\mathrm{gr}}(G, S\gamma) \xrightarrow{\ \sim\ } \underline{\mathrm{Ext}}^{i+n}_{k\text{-}\mathrm{gr}}(T^{*}G, \gamma)
\]
\marginal{$S\gamma = (D \Delta \gamma_{K})(1)$}
\note{la note encadrée en haut à gauche, reliée par une flèche au premier
$S\gamma$, donne la définition de $S\gamma$ ; le $D$ en est une lecture
douteuse. Ses $\underline{\mathrm{Ext}}$ sont soulignés deux fois, d'un trait
droit et d'un trait ondulé ; on ne rend que le premier. L'exposant $i+n$ est
surchargé, de même le second argument du membre de droite (un $\gamma$ sur une
autre lettre) ; sur $T^{*}$ l'astérisque est surchargé.}
isom.\ si $G$ \emph{affine} sans composante $\mathbb{G}_{a}$
\[
\underline{\mathrm{Ext}}^{i}_{K/k}(G, S\gamma) \to \underline{\mathrm{Ext}}^{2i}(T^{*}G, \gamma)
\]
\note{l'exposant « $2i$ » est surchargé ; la lecture n'en est pas sûre.}
isom.
\[
\underline{\mathrm{Ext}}^{i}_{K/k}(G, S\gamma) \to \underline{\mathrm{Ext}}^{i}
\]
\note{la ligne s'arrête là. Plus bas, à gauche, un symbole biffé (\struck{\ill{}}) indicé $K/k$, puis, seul, « $\mathrm{R}^{*}\!\int_{K/k}$ ».}
\begin{enumerate}
\item[1)] $K$-\emph{groupes algébriques}. $\mathrm{Ext}$,
$\underline{\mathrm{Ext}}$, $H^{p}$ \emph{de} $K$-\emph{groupes algébriques}
\marginal{nuls pour $i \geqslant 2$}
\note{sous $\mathrm{Ext}$, $\underline{\mathrm{Ext}}$ et $H^{p}$, trois traits
descendent vers trois gloses serrées : « \uncertain{nature}
\uncertain{arithmétique} \uncertain{i.e.} \uncertain{globale} »,
« \uncertain{nature} \uncertain{géométrique} », « \uncertain{nature}
\uncertain{arithmétique} \uncertain{i.e.}\ \uncertain{globale} » ; « nuls pour
$i \geqslant 2$ » est relié par un trait à $\underline{\mathrm{Ext}}$.}
\item[2)] Dualité de Cartier. Extension aux \ill{}$\mathrm{gr}_{K}$
\item[3)] Cohomologies algébriques, leurs $\mathrm{Ext}$,
$\underline{\mathrm{Ext}}$, $H^{p}$
\item[4)] \emph{Dualité de \uncertain{Serre-Rosenlicht}}
\item[5)] Le foncteur $\underline{\Gamma}_{K/k}$ \struck{\ill{}} (le Néron-\ill{})
\marginal{nul pour $i > 3$}
\item[6)] Les foncteurs \struck{\ill{}} $\underline{\mathrm{Ext}}^{*}_{K/k}(F, G)$,
$\underline{H}^{*}_{K/k}(G)$
\[
E^{*}_{K/k}(F, G) \qquad \mathbf{E}_{K/k}(G)
\]
\item[7)] L'hom.\ fondamental.
\item[8)] Le th.\ de dualité.
\end{enumerate}
\note{plan de travail numéroté 1)--8), au crayon ; les items 3), 5) et 8)
portent un double trait oblique dans la marge gauche, sans doute une coche. Au
point 2), le symbole devant $\mathrm{gr}_{K}$ est une grande lettre surchargée
(un $\mathcal{K}$ ? un « Ind » ?). Au point 4), le titre est souligné d'un trait qui
remonte encadrer la parenthèse de 5), « (le Néron-\ill{}) », dont le second nom
est surchargé ; le « nul pour $i > 3$ » y est relié. Au point 6), un mot est
biffé avant $\underline{\mathrm{Ext}}^{*}_{K/k}$, et $\underline{H}^{*}_{K/k}$
est souligné d'un trait ondulé ; $E^{*}_{K/k}(F, G)$ et
$\mathbf{E}_{K/k}(G)$ sont écrits dessous comme notations abrégées.}

\end{document}
